\documentclass[aspectratio=169,11pt]{beamer}

\usetheme{Madrid}
\usecolortheme{default}
\usefonttheme{professionalfonts}
\setbeamertemplate{navigation symbols}{}
\setbeamertemplate{footline}[frame number]

\usepackage{amsmath, amssymb, booktabs, graphicx}

\newcommand{\E}{\mathbb{E}}
\newcommand{\1}{\mathbf{1}}
\newcommand{\Prb}{\mathbb{P}}

\title{When and Why Structure? Static and Dynamic Decision Problems}
\subtitle{Empirical Methods --- Week 6}
\author{Teaching deck draft}
\date{}

\begin{document}

\begin{frame}
  \titlepage
\end{frame}

\begin{frame}{Week 6 question}
\begin{itemize}
    \item When does an applied question require an economic model rather than only a credible comparison?
    \item What is observed, what is latent, and what variation disciplines the latent objects?
    \item When do expectations, dynamics, welfare, or equilibrium response change the research design?
    \item How should a reader interpret a structural counterfactual?
\end{itemize}

\medskip
Structure is a research choice, not a rival ideology to design-based work.
\end{frame}

\begin{frame}{Why structure?}
\begin{itemize}
    \item Latent primitives: preferences, costs, technologies, beliefs, skills, frictions.
    \item Expectations: decisions depend on future wages, offers, policy rules, or risks.
    \item Dynamic incentives: actions today change tomorrow's state.
    \item Counterfactuals: the policy of interest may be outside observed support.
    \item Welfare: utility, equivalent variation, and insurance value are not directly observed.
    \item Equilibrium response: prices, queues, vacancies, and market composition may adjust.
\end{itemize}

\medskip
The model should buy an object that the data cannot reveal transparently without structure.
\end{frame}

\begin{frame}{Static vs dynamic problems}
\small
\begin{tabular}{p{0.18\linewidth} p{0.26\linewidth} p{0.25\linewidth} p{0.21\linewidth}}
\toprule
Type & Example & Key latent object & Main risk \\
\midrule
Static choice & job or school menu & contemporaneous utility tradeoffs & wrong menu or shock structure \\
Static labor supply & hours under tax schedule & preferences over consumption and leisure & functional-form dependence \\
Dynamic schooling & school, work, occupation & continuation values and skill accumulation & state-space misspecification \\
Dynamic search & accept offer or search & offer distribution and reservation value & expectations and duration dependence \\
Lifecycle labor supply & work, saving, retirement & risk, insurance, intertemporal preferences & extrapolation and heterogeneity \\
\bottomrule
\end{tabular}
\end{frame}

\begin{frame}{Observed and latent objects}
\small
\begin{tabular}{p{0.25\linewidth} p{0.29\linewidth} p{0.34\linewidth}}
\toprule
Observed in data & Latent in model & Identification question \\
\midrule
choices and transitions & preferences and switching costs & which choices shift when states or prices change? \\
wages and employment & skill prices and productivity & which histories discipline returns and shocks? \\
policy rules and prices & perceived incentives and beliefs & do agents know and respond to the rule? \\
state histories & transition laws and option values & which state changes predict future behavior? \\
outcomes & welfare and compensating variation & how does utility map to observed behavior? \\
\bottomrule
\end{tabular}

\medskip
A structural paper should make this map visible before the main counterfactual.
\end{frame}

\begin{frame}{Static latent utility}
\[
U_{ij}
=
u_j(X_i,Z_{ij};\theta)
+ \alpha_i q_{ij}
+ \xi_j
+ \varepsilon_{ij}
\]

\[
d_{ij}
=
\1\{U_{ij}\ge U_{ik}\ \text{for all }k\in\mathcal J_i\}
\]

\[
\Prb(d_{ij}=1\mid X_i,Z_i;\theta)
=
\frac{\exp(V_{ij})}
{\sum_{k\in\mathcal J_i}\exp(V_{ik})}
\]

\begin{itemize}
    \item Observed: choice, menu, attributes, prices or wages.
    \item Latent: utility, nonpecuniary value, unobserved quality, shocks.
    \item The shock assumption turns latent utility into choice probabilities.
\end{itemize}
\end{frame}

\begin{frame}{Dynamic Bellman logic}
\[
V(s_t)
=
\max_{a_t\in\mathcal A(s_t)}
\left\{
u(a_t,s_t;\theta)
+
\delta\E_{\theta}
\left[
V(s_{t+1})\mid s_t,a_t
\right]
\right\}
\]

\[
v(a,s;\theta)
=
u(a,s;\theta)
+
\delta\sum_{s'}V(s';\theta)F(s'\mid s,a;\theta)
\]

\begin{itemize}
    \item States summarize payoff-relevant history.
    \item Actions change future states and future options.
    \item Continuation values are the object reduced-form comparisons often leave latent.
\end{itemize}
\end{frame}

\begin{frame}{Continue vs stop}
\[
v_C(s;\theta)
=
u_C(s;\theta)
+
\delta\sum_{s'}V(s';\theta)F_C(s'\mid s;\theta)
\]

\[
v_S(s;\theta)
=
u_S(s;\theta)
+
\delta\sum_{s'}V(s';\theta)F_S(s'\mid s;\theta)
\]

\[
\Prb(a=S\mid s;\theta)
=
\frac{\exp(v_S(s;\theta))}
{\exp(v_S(s;\theta))+\exp(v_C(s;\theta))}
\]

\begin{itemize}
    \item Replacement: keep operating or replace an engine.
    \item Search: accept a job or continue searching.
    \item Schooling: stay in school or work now.
\end{itemize}
\end{frame}

\begin{frame}{Estimation strategies}
\small
\textbf{Likelihood}
\[
\widehat\theta
=
\arg\max_{\theta}
\sum_{i,t}
\log \Prb_{\theta}(a_{it}\mid s_{it},p_{it})
+
\log f_{\theta}(y_{it},s_{i,t+1}\mid a_{it},s_{it},p_{it})
\]

\textbf{Moments or simulation}
\[
\widehat\theta
=
\arg\min_{\theta}
\left[\widehat m^{data}-m^{model}(\theta)\right]'
W
\left[\widehat m^{data}-m^{model}(\theta)\right]
\]

\begin{itemize}
    \item Fit should report targeted moments and non-target validation.
    \item Identification asks which variation disciplines each primitive.
\end{itemize}
\end{frame}

\begin{frame}{Policy counterfactuals and welfare}
\[
\widehat\theta
\Longrightarrow
\left\{
\sigma_{\widehat\theta}(a\mid s;p),
F_{\widehat\theta}(s'\mid s,a;p),
Y_{\widehat\theta}(p),
W_{\widehat\theta}(p)
\right\}
\]

\[
\Delta Y(p_1,p_0)
=
\E_{\widehat\theta}[Y_i(p_1)-Y_i(p_0)],
\qquad
\Delta W(p_1,p_0)
=
\E_{\widehat\theta}[W_i(p_1)-W_i(p_0)]
\]

\begin{itemize}
    \item The policy changes rules, prices, constraints, or menus.
    \item The model holds some primitives fixed and recomputes behavior.
    \item Credibility falls as the counterfactual moves away from identifying support.
\end{itemize}
\end{frame}

\begin{frame}{Theory to applied papers}
\small
\begin{tabular}{p{0.19\linewidth} p{0.25\linewidth} p{0.25\linewidth} p{0.21\linewidth}}
\toprule
Paper & Question & What is latent? & Model contribution \\
\midrule
Rust & replace now or later? & replacement costs and value of waiting & dynamic replacement counterfactuals \\
Hotz-Miller & estimate dynamic choice? & value differences & CCP-based estimation logic \\
Keane-Wolpin & schooling, work, career? & skills, preferences, expectations & dynamic career paths \\
Todd-Wolpin & schooling subsidy? & household behavioral rules & validation and policy simulation \\
Low-Meghir-Pistaferri & lifecycle risk? & insurance and risk processes & welfare under wage and employment risk \\
Adda-Dustmann-Stevens & child career costs? & dynamic timing and persistence & long-run career counterfactuals \\
\bottomrule
\end{tabular}
\end{frame}

\begin{frame}{Credibility threats}
\begin{itemize}
    \item Functional form: results driven by parametric utility, shocks, or transitions.
    \item Expectations: agents may not know or believe the process assumed by the model.
    \item Latent heterogeneity: flexible types can absorb misspecification.
    \item Off-target fit: the model matches targeted moments but misses relevant behavior.
    \item Identification by assumption: key primitives come from normalization more than data.
    \item Extrapolation: the counterfactual changes support, menus, or equilibrium conditions.
\end{itemize}

\medskip
Structural discipline means reporting these limits, not hiding them.
\end{frame}

\begin{frame}{Research Lab design}
\begin{columns}[T,onlytextwidth]
\begin{column}{0.32\linewidth}
\textbf{Reproduce}
\begin{itemize}
    \item Rust-style replacement
    \item Mileage states
    \item Replacement CCPs
    \item Dynamic logit likelihood
\end{itemize}
\end{column}
\begin{column}{0.32\linewidth}
\textbf{Diagnose}
\begin{itemize}
    \item Static vs dynamic fit
    \item State support
    \item Parameter sensitivity
    \item Replacement-cost counterfactual
\end{itemize}
\end{column}
\begin{column}{0.32\linewidth}
\textbf{Transfer}
\begin{itemize}
    \item Career-choice panel
    \item School/work states
    \item Tuition subsidy
    \item Counterfactual memo
\end{itemize}
\end{column}
\end{columns}
\end{frame}

\begin{frame}{Takeaways}
\begin{itemize}
    \item Use structure when the research object is latent, dynamic, welfare-relevant, or outside observed support.
    \item The model must identify what is observed, what is latent, and what data variation disciplines the mapping.
    \item Bellman logic is the language of dynamic incentives and continuation values.
    \item Estimation is credible only when fit, validation, and sensitivity are reported transparently.
    \item Counterfactuals should be interpreted through the distance between the new policy and the evidence that identified the primitives.
\end{itemize}
\end{frame}

\end{document}
